\section{Modelling Pipeline and Model Selection}
\label{sec:pipeline}

All models share one common pipeline structure. A model comparison is 
only meaningful when every model receives the same rows, the same target, 
the same split, and the same evaluation metrics. The repository therefore 
separates shared utilities from model-specific scripts.

\begin{table}[htbp]
\caption{Pipeline stages and their purpose.}
\label{tab:pipeline}
\centering
\scriptsize
\begin{tabularx}{\columnwidth}{p{0.38\columnwidth} X}
\toprule
\textbf{Stage} & \textbf{Purpose} \\
\midrule
Raw data loading & 
Read trip, station, and weather files. \\
\addlinespace
Station-day aggregation & 
Aggregate rentals and merge weather by date. \\
\addlinespace
Feature reduction & 
Remove leakage columns and redundant predictors. \\
\addlinespace
Lag/no-lag variants & 
Create two comparable datasets with identical row keys. \\
\addlinespace
One-hot encoding & 
Encode 20 station identities as 19 one-hot columns. \\
\addlinespace
Model-specific scaling & 
Standardize inputs only for models that require it. \\
\addlinespace
Chronological split & 
70\% training, 15\% validation, 15\% test. \\
\addlinespace
Training and tuning & 
Fit all models, then optimize hyperparameters on 
the validation set. \\
\addlinespace
Evaluation & 
Save metrics, predictions, rankings, and plots. \\
\bottomrule
\end{tabularx}
\end{table}

\subsection{Chronological Split and Scaling}

A random split would mix past and future observations, making the task 
easier than real forecasting. The split is therefore strictly 
chronological: 70\% training, 15\% validation, 15\% test. The 
validation set is used for all tuning decisions. The test set is 
only used once for final reporting.

Scaling is applied only where the model needs it. KNN uses distances, 
regularized linear models use coefficient penalties, and neural networks 
use gradient-based optimization -- all sensitive to input scale. 
Decision Tree, Random Forest, and Gradient Boosting split data by 
thresholds and are not affected by scale. We did not apply PCA because 
every feature has a clear physical meaning and replacing them with 
abstract components would make results harder to explain.

\subsection{Model Selection}

The nine models were chosen to cover increasing complexity and different 
learning principles from the course. This is more informative than 
testing many similar variants of one algorithm.

\begin{table}[htbp]
\caption{Models, input variant, scaling, and reason for inclusion.}
\label{tab:models}
\centering
\scriptsize
\begin{tabularx}{\columnwidth}{p{0.26\columnwidth} c p{0.08\columnwidth} X}
\toprule
\textbf{Model} & \textbf{Lag} & \textbf{Scale} & \textbf{Reason} \\
\midrule
Dummy Regressor & no & no & 
Minimum baseline; all real models must beat it. \\
\addlinespace
Linear Regression & both & yes & 
Tests whether features contain a roughly linear signal. \\
\addlinespace
Ridge & both & yes & 
L2 penalty checks whether shrinkage helps. \\
\addlinespace
Lasso & both & yes & 
L1 penalty adds automatic feature selection. \\
\addlinespace
Decision Tree & both & no & 
Tests nonlinear threshold rules; prone to overfitting. \\
\addlinespace
KNN & both & yes & 
Tests whether similar station-days have similar demand. \\
\addlinespace
Random Forest & both & no & 
Bagging reduces single-tree variance. \\
\addlinespace
Gradient Boosting & both & no & 
Sequential trees target remaining prediction errors \cite{friedman2001}. \\
\addlinespace
Neural Network & both & yes & 
Flexible nonlinear reference to compare against tree ensembles. \\
\bottomrule
\end{tabularx}
\end{table}
